\documentclass[11pt]{article}

\usepackage[margin=1in]{geometry}
\usepackage{amsmath,amssymb,graphicx,multicol,array,enumerate,textcomp}
\usepackage{tikz}
\usepackage{pgfplots}
\pgfplotsset{compat=1.18}
\usetikzlibrary{arrows.meta,calc}

\newcommand{\usd}{\text{\$}}

\newenvironment{problem}[2][Problem]{\begin{trivlist}
\item[\hskip \labelsep {\bfseries #1}\hskip \labelsep {\bfseries #2.}]
\leavevmode\par
}{\end{trivlist}}

\begin{document}

\title{Problem Set 5:}
\author{Joe White\\ECO 442}
\date{\today}
\maketitle

% =========================================================
\begin{problem}{1}

\vspace{0.5em}
\begin{enumerate}[(a)]
\item \textbf{Decrease prices in Costa Rica in the long run. Should you raise or lower the money supply?}

In the long run, money market equilibrium is
\[
MS_C = MD_C
\quad\Longleftrightarrow\quad
MS_C = P_C \cdot L(R_C, Y_C),
\]
or
\[
\frac{MS_C}{P_C} = L(R_C, Y_C).
\]
In the long run, $P_C$ can change, so a if we want lower prices $P_C$ in the long run, we need a lower nominal money supply:
\[
MS_C \downarrow \ \Rightarrow\ P_C \downarrow \quad
\]

\item \textbf{PPP for the col\'on/USD exchange rate. What happens when $P_C$ falls?}

Absolute PPP implies
\[
E_{C/\usd} \;=\; \frac{P_C}{P_{\usd}}.
\]
Holding the U.S. price level $P_{\usd}$ fixed, a fall in Costa Rica's price level implies
\[
P_C \downarrow \ \Rightarrow\ E_{C/\usd} \downarrow.
\]

\[
E_{C/\usd}\downarrow \quad\Longleftrightarrow\quad \text{the col\'on appreciates vs.\ the dollar}
\]

\item \textbf{Affordability for a Costa Rican consumer buying a car produced in Costa Rica (long run).}

$MS_C\downarrow$, which in the long run implies $P_C\downarrow$

\begin{enumerate}[(1)]
\item \textbf{Price of the car in Costa Rica (in col\'ones):} Its nominal price moves with the price level in the long run, so
\[
P^{car}_C \downarrow \quad \text{(proportional to } P_C\downarrow\text{).}
\]

\item \textbf{Financing the Costa Rican consumer can receive:}
In the long run, real interest rates are pinned down by real factors. However, lower long-run inflation implies lower nominal rates (Fisher effect). So nominal borrowing rates tend to be lower, but the real cost of borrowing is unchanged.

\item \textbf{Nominal income of the Costa Rican consumer:}
With money neutrality, nominal income tends to move proportionally with the price level:
\[
\text{Nominal income}_C \downarrow \quad \text{(roughly with } P_C\downarrow\text{).}
\]

\item \textbf{Real income of the Costa Rican consumer:}
Real income is a real variable and is not changed by a permanent level shift in money in the long run:
\[
\text{Real income}_C \approx \frac{\text{Nominal income}_C}{P_C} \ \text{is unchanged.}
\]
\end{enumerate}

The nominal price of the car falls, but nominal income also falls proportionally, leaving real purchasing power roughly unchanged. So in real terms, the affordability of a Costa Rican made car for a Costa Rican consumer is unchanged in the long run.

\item \textbf{Affordability for a U.S. consumer buying a car produced in Costa Rica (long run).}

\begin{enumerate}[(1)]
\item \textbf{Price of the car in the U.S. (in dollars):}

\[
P^{car}_{\usd} = \frac{P^{car}_C}{E_{C/\usd}}.
\]
In the long run, $P^{car}_C$ falls with $P_C$, and $E_{C/\usd} = P_C/P_{\usd}$ also falls proportionally when $P_C$ falls.
So the ratio is unchanged:
\[
\frac{P^{car}_C \downarrow}{E_{C/\usd}\downarrow} = \text{no change in } P^{car}_{\usd}.
\]


\item \textbf{Financing the U.S. consumer can receive:}
This policy is in Costa Rica; in the long run it does not directly change U.S. real interest rates or U.S. credit conditions.

\item \textbf{Nominal income of the U.S. consumer:} unchanged.

\item \textbf{Real income of the U.S. consumer:} unchanged.
\end{enumerate}

\end{enumerate}
\end{problem}

% =========================================================
\begin{problem}{2}

\[
g_{JP}=1\%, \quad g_{KR}=6\%, \qquad
\mu_{JP}=4\%, \quad \mu_{KR}=7\%.
\]

\begin{enumerate}[(a)]
\item \textbf{Inflation rate in each country.}

Using the monetary model in growth rates:
\[
\pi = \mu - g.
\]
So
\[
\pi_{JP} = 4\% - 1\% = 3\%,
\qquad
\pi_{KR} = 7\% - 6\% = 1\%.
\]

\item \textbf{Which currency is depreciating and at what annual rate?}

Let $E_{W/\yen}$ be the exchange rate in won per yen. Then
\[
\%\Delta E_{W/\yen} = \pi_{KR} - \pi_{JP} = 1\% - 3\% = -2\%.
\]
So $E_{W/\yen}$ falls by $2\%$ per year, meaning:
\[
E_{W/\yen}\downarrow \ \Rightarrow\ \text{the won appreciates and the yen depreciates at } 2\%\text{/year.}
\]

\item \textbf{Korea wants inflation fixed at 2\% per year. What money growth should they target?
Then explain effects on (i) real rates, (ii) nominal rates, (iii) the won/yen exchange rate (long run).}

Target $\pi_{KR}=2\%$ and keep $g_{KR}=6\%$. From $\pi=\mu-g$:
\[
2\% = \mu_{KR} - 6\%
\quad\Longrightarrow\quad
\mu_{KR} = 8\%.
\]
So the Bank of Korea should raise money growth from $7\%$ to $8\%$.

\begin{enumerate}[(i)]
\item \textbf{Real interest rates in Korea:} In the long run, real interest rates are determined by real factors and do not change systematically with a change in long-run inflation. So real rates are unchanged.

\item \textbf{Nominal interest rates in Korea:} By the Fisher effect, nominal rates move one-for-one with inflation in the long run. Inflation rises from $1\%$ to $2\%$, so nominal interest rates rise by about $1$ percentage point.

\item \textbf{Exchange rate between the won and yen (long run):}
Japan stays at $\pi_{JP}=3\%$. Korea becomes $\pi_{KR}=2\%$. Relative PPP implies
\[
\%\Delta E_{W/\yen} = \pi_{KR} - \pi_{JP} = 2\% - 3\% = -1\%.
\]
So the won still appreciates (and the yen still depreciates), but now only at 1\% per year instead of 2\% per year.
\end{enumerate}

\end{enumerate}
\end{problem}

\end{document}